\documentclass[11pt]{letter}
\usepackage[a4paper,margin=1in]{geometry}
\usepackage{hyperref}
\setlength{\parindent}{0pt}
\setlength{\parskip}{8pt}
\signature{Hikaru Wakaura}
\address{QIRI, Quantum Integrated Research Institute Inc.\\ Tokyo 107-0061, Japan\\ \texttt{h.wakaura@qiri.co.jp}}
\begin{document}
\begin{letter}{The Editors\\ \textit{Physical Chemistry Chemical Physics}}

\opening{Dear Editors,}

On behalf of my co-author, I am submitting the manuscript \textbf{``What a
radical-pair quantum reservoir would require: readable capacity, no quantum
advantage, and a nuclear-register criterion''} for consideration as a research
article in \textit{PCCP}.

This is a re-submission, and I should say plainly what happened. It supersedes
manuscript \textbf{CP-ART-06-2026-002404}, which we asked to withdraw during review
after discovering ourselves that one of its central numerical results was an
artifact of our simulation protocol rather than a property of the spin chemistry.
The present manuscript is the corrected and substantially extended analysis. Its
conclusions differ from the withdrawn version in three places, and in each case the
corrected answer is the more interesting one. We are grateful for the opportunity to
put the corrected work before the journal.

The question is whether the spin chemistry of a cryptochrome-type radical pair can
\emph{process} information---posed, as it must be for a warm and noisy system, as
reservoir computing rather than gate-model computation. We solve the Liouville-space
Haberkorn master equation with electronic decoherence and, crucially, drive the
reservoir the way a real turnover does: the nuclei are retained and a new
spin-correlated pair is born each cycle, its singlet character encoding the input.
Resetting a single electron instead---the natural abstraction, and the one used in
the withdrawn version---destroys the electron--electron correlation and, for a
separated pair, pins the singlet probability at exactly $1/4$ regardless of input.
That is the artifact, and correcting it changes the physics.

Three results follow. \textbf{First}, the reservoir is readable by ordinary
chemistry. The time-resolved yield of one product pool gives an out-of-sample
information-processing capacity of $2.0$, and the nuclear polarisation the product
carries away raises it to $4.7$, with no spectroscopic access of any kind. Readout
is not the bottleneck, as we had previously concluded. The mechanism behind that
number is itself worth reporting to a physical-chemistry readership: because either
electron of a newly born pair is maximally mixed for every input, spin-selective
recombination is what \emph{writes} the input into the nuclear register. In a
separated pair, CIDNP is not an optional additional channel; it is the write
mechanism.

\textbf{Second}, there is no quantum advantage. A classical echo-state network given
the same number of readout features matches or exceeds the five-spin reservoir. We
report this rather than argue around it.

\textbf{Third}, the timescale objection usually raised against this system does not
survive. It rests on identifying the reservoir's clock with the microsecond
radical-pair lifetime, but the memory is carried by the nuclear register, which is
not consumed by recombination and departs in the diamagnetic product. Inserting a
physiological $1$--$100$~ms turnover leaves the memory capacity unchanged to within
$0.2\%$ and places the horizon at $3$--$280$~ms, inside the neural band. What
replaces the timescale argument is a sharper and falsifiable chemical criterion: the
\emph{same} nuclear register must be reused from one turnover to the next, failing
which the capacity collapses to a memoryless read-back. Unlike a comparison of two
timescales, which requires no quantum mechanics at all, this criterion can only be
evaluated by the kind of spin-dynamical analysis performed here---and it is directly
testable in synthetic donor--bridge--acceptor radical pairs by cycling the system and
asking whether the yield at one turnover depends on the input at the previous one.

The work is entirely computational. All code and numerical outputs are openly
released in the \texttt{qbscreen} package, every capacity is reported out of sample
with its null floor and its convergence in sample length, and the central claims are
covered by an automated test suite (28 tests, all passing). A companion study on the
corresponding weak-field magnetosensing analysis is under consideration elsewhere;
the two are physically independent and we are glad to share it. This manuscript is
not under consideration at any other journal. The authors declare no competing
interests and the study received no specific funding.

Thank you for considering this re-submission.

\closing{Yours sincerely,}

\end{letter}
\end{document}
